% Figura V5: construção da solução parcial (3) -> (3,1) -> (3,1,4) -> (3,1,4,2)
\begin{tikzpicture}[
    scale=0.9, transform shape,
    >=Stealth,
    cel/.style={draw=btNeutro, minimum size=6.5mm, font=\small, inner sep=0pt},
    cheia/.style={cel, draw=btValido, fill=btValido!18, text=btTexto},
    vazia/.style={cel, fill=btFundo},
    rotulo/.style={font=\scriptsize, text=btNeutro},
  ]
  % Cada quadro: 4 células; as preenchidas mostram os números escolhidos.
  % Frame 1: (3)
  \begin{scope}[shift={(0,0)}]
    \node[cheia] at (0,0) {3};
    \node[vazia] at (0.7,0) {}; \node[vazia] at (1.4,0) {}; \node[vazia] at (2.1,0) {};
    \node[rotulo] at (1.05,-0.65) {Etapa 1};
  \end{scope}
  % Frame 2: (3,1)
  \begin{scope}[shift={(3.6,0)}]
    \node[cheia] at (0,0) {3}; \node[cheia] at (0.7,0) {1};
    \node[vazia] at (1.4,0) {}; \node[vazia] at (2.1,0) {};
    \node[rotulo] at (1.05,-0.65) {Etapa 2};
  \end{scope}
  % Frame 3: (3,1,4)
  \begin{scope}[shift={(7.2,0)}]
    \node[cheia] at (0,0) {3}; \node[cheia] at (0.7,0) {1}; \node[cheia] at (1.4,0) {4};
    \node[vazia] at (2.1,0) {};
    \node[rotulo] at (1.05,-0.65) {Etapa 3};
  \end{scope}
  % Frame 4: (3,1,4,2)
  \begin{scope}[shift={(10.8,0)}]
    \node[cheia] at (0,0) {3}; \node[cheia] at (0.7,0) {1};
    \node[cheia] at (1.4,0) {4}; \node[cheia] at (2.1,0) {2};
    \node[rotulo] at (1.05,-0.65) {Etapa 4};
  \end{scope}

  % setas entre os quadros
  \foreach \x in {2.7, 6.3, 9.9}{
    \draw[->, thick, draw=btNeutro] (\x,0) -- (\x+0.5,0);
  }
\end{tikzpicture}
